\section{Related work}\label{sec:related}

\subsection{Reinforcement learning for bus holding control}
\label{subsec:rl_holding}

Research on real-time bus holding spans analytical control rules and learned policies. Early analytical rules remain important reference points. Daganzo~\citep{daganzo2009headway} derived a headway-based control law for eliminating bunching. Xuan et al.~\citep{xuan2011dynamic} proposed dynamic holding strategies that balance schedule reliability and regular headways through optimal linear-control analysis. We include corresponding classical baselines to test whether learning adds value beyond interpretable holding rules.

RL studies learn holding decisions directly from simulated operations. Alesiani and Gkiotsalitis~\citep{alesiani2018reinforcement} formulated high-frequency bus holding as a value-based RL problem. Wang and Sun~\citep{wang2020real} introduced a multi-agent deep RL controller in which each bus acts as an agent under shared training. Zhang et al.~\citep{zhang2026single} developed a single-agent robust deep RL framework for bidirectional timetabled bus lines, showing that categorical identifiers help a shared policy operate across heterogeneous vehicles, stations, and trips.

Subsequent work addressed asynchronous control and robustness. Wang and Sun~\citep{wang2021asynchronous} developed an asynchronous multi-agent reinforcement learning (MARL) framework with counterfactual credit assignment, while Wang and Sun~\citep{wang2023robust} combined distributional MARL with meta-learning to handle service perturbations. Multi-line and multi-strategy settings include shared-corridor multi-objective MARL \citep{wang2023multiobjective}, cooperative holding with stop-skipping \citep{rodriguez2023cooperative}, hierarchical MARL that combines holding and speed control \citep{yu2024hierarchical}, and graph-aware deep RL for multi-line busy corridors \citep{li2026graphaware}.

Other studies focus on RL design choices. He et al.~\citep{he2022dynamic} used Q-learning with multistage look-ahead for adaptive holding. Xu et al.~\citep{xu2025systematic} compared state, action, and agent configurations for RL-based bunching mitigation. Yu et al.~\citep{yu2025llm} used language-model-assisted reward generation for generic bus holding policies. Zhang and Zheng~\citep{resac2026} proposed robust ensemble soft actor-critic (RE-SAC), an ensemble-critic variant of soft actor-critic for bus fleet control that separates aleatoric and epistemic risk signals. We adopt this critic structure at both levels of TransitDuet.

Many RL-based holding controllers treat the timetable, target headway, or dispatch plan as an exogenous input. They can react to realised bunching but do not use systematic execution evidence---such as repeated lateness, excessive onboard delay, or recurrent over-holding---to revise future dispatch times. TransitDuet addresses this limitation by coupling a holding controller with a higher-level rolling timetable policy that feeds execution outcomes back into dispatch planning.

\subsection{Bus timetable optimisation}
\label{subsec:timetable_opt}

Timetable design has traditionally been approached through mathematical programming and metaheuristic search. Ibarra-Rojas et al.~\citep{ibarra2015planning} formulated a mixed-integer programme that jointly determines departure times and vehicle assignments to minimise a weighted combination of passenger waiting time and operator cost. They solved it by branch-and-bound with valid inequalities. Gkiotsalitis and Cats~\citep{gkiotsalitis2021bus} proposed a bi-level model in which the upper level selects headways to minimise total passenger cost and the lower level solves a vehicle-scheduling problem using adaptive large neighbourhood search. Frequency-setting variants, such as the work of Cats and Gl{\"u}ck~\citep{cats2019frequency}, optimise the number of departures per period rather than individual departure times. They typically rely on assignment models to estimate passenger flows under each candidate frequency vector.

Many timetable formulations rely on static demand forecasts and assumed travel-time conditions. They are optimised offline against expected conditions and then handed off to operations without a feedback loop. Under traffic incidents, demand surges, or driver variability, the timetable may become infeasible, leaving holding controllers to absorb the mismatch. The combinatorial structure can also limit scalability as the number of trips and stops increases. TransitDuet instead uses a closed-loop approach in which the upper policy maps aggregate state features to planned-headway adjustments and receives lower-level execution feedback through the cross-level coupling described in Section~\ref{sec:tap}.

\subsection{Hierarchical and multi-timescale reinforcement learning}
\label{subsec:hierarchical_rl}

Hierarchical reinforcement learning decomposes complex tasks across multiple temporal scales. The options framework of Sutton et al.~\citep{sutton1999between} introduced temporally abstract actions, each consisting of an initiation set, an intra-option policy, and a termination condition. It enables planning over multiple timescales within a single Markov decision process (MDP). Feudal networks \citep{vezhnevets2017feudal} operationalize this idea with a manager that sets subgoals in a learned latent space and a worker that executes primitive actions to reach them. HIRO \citep{nachum2018data} introduced off-policy corrections that relabel high-level goals in hindsight, making hierarchical training practical in continuous-control domains. Lee et al.~\citep{lee2020reinforcement} and Metelli et al.~\citep{metelli2020control} studied multi-timescale settings more broadly, including interactions among agents operating at different decision frequencies.

These methods provide useful abstractions for temporal hierarchy, but many assume synchronous execution or a shared state space across levels. In transit operations, the timetable agent acts at dispatch events, whereas the holding agent acts at every stop visit. The two levels therefore observe different state representations at non-aligned decision epochs. TransitDuet handles this asynchrony with a dynamic timetable channel: the upper selects planned-headway adjustments, and the lower executes the resulting timetable through holding.

\subsection{Constrained reinforcement learning}
\label{subsec:constrained_rl}

Many practical RL deployments must satisfy constraints beyond reward maximisation. Lagrangian relaxation is a widely used approach. Tessler et al.~\citep{tessler2019reward} proposed reward-constrained policy optimisation, which augments the policy gradient with dual variables corresponding to cost constraints and performs primal-dual updates to steer the policy toward feasibility. Stooke et al.~\citep{stooke2020responsive} introduced proportional-integral-derivative (PID)-based Lagrange multiplier updates to reduce oscillation and accelerate convergence. These methods handle soft constraints, for which occasional small violations are tolerable, but do not provide worst-case guarantees.

For settings that demand hard constraint satisfaction, Miryoosefi et al.~\citep{miryoosefi2019reinforcement} introduced ApproPO. It formulates the constrained problem as online convex optimisation over occupancy measures and applies a measurement-based projection oracle to enforce feasibility at every iterate. The $\theta$-OGD variant provides a practical implementation that projects policy parameters onto the feasible set using online gradient descent with approximate projections. In the multi-agent domain, Pan et al.~\citep{pan2024roboduet} proposed RoboDuet for humanoid locomotion. Its upper-body controller and lower-body locomotor exchange cross-advantage terms without centralised control.

TransitDuet combines elements from both lines of work to address the constraint structure of transit operations. Fleet-size feasibility is handled as a soft constraint through the $\theta$-OGD measurement-projection scheme of Miryoosefi et al.~\citep{miryoosefi2019reinforcement}. This scheme adaptively reweights an overshoot penalty in the upper-level reward. Per-episode fluctuations under elastic fleet sampling and demand noise can still produce small, bounded overshoot. Headway uniformity is handled as a soft constraint through a Lagrangian penalty with adaptive multiplier updates in the lower-level holding policy. The dynamic timetable channel in Section~\ref{sec:tap} manages cross-level coordination.
